\chapter{Introducere}
\label{chap:introducere}

\section{Contextul Cercetării și Motivare}
Îmbătrânirea demografică accelerată la nivel global și creșterea ponderii persoanelor vârstnice care locuiesc singure impun dezvoltarea unor soluții tehnologice fiabile și non-invazive de asistență medicală și monitorizare a siguranței la domiciliu. Conform rapoartelor Organizației Mondiale a Sănătății (OMS), căderile accidentale reprezintă a doua cauză principală a deceselor provocate de traumatisme neintenționate la nivel mondial, peste 37.3 milioane de căderi necesitând îngrijiri medicale specializate în fiecare an. În mod deosebit, fenomenul cunoscut în literatura de specialitate sub denumirea de ,,long lie'' (rămânerea la sol timp de peste o oră după producerea evenimentului din cauza imposibilității de auto-redresare) crește exponențial riscul de deshidratare severă, hipotermie, rabdomioliză și deces precoce.

Soluțiile clasice existente pe piață prezintă limitări structurale majore care le reduc eficacitatea în scenarii reale de viață:
\begin{enumerate}
    \item \textbf{Sistemele bazate pe senzori purtabili (wearables):} Brățările inteligente, pandantivele de urgență cu accelerometre și giroscoape sau ceasurile inteligente necesită purtare continuă și reîncărcare periodică. Studiile clinice indică o rată de complianță sub 45\% în rândul vârstnicilor, mulți uitând să le poarte sau refuzându-le pe timpul nopții ori în cabina de duș—zone cu cel mai ridicat risc de alunecare.
    \item \textbf{Sistemele bazate pe viziune computerizată (camere video optice sau infraroșu):} Deși oferă o precizie ridicată a recunoașterii tiparelor posturale, camerele optice sunt inacceptabile în spațiile intime (dormitor, baie) din considerente stricte de viață privată. În plus, funcționarea lor este afectată de condițiile de iluminare ambientală scăzută și de obstrucțiile vizuale directe generate de mobilier.
    \item \textbf{Sistemele acustice sau pardoselile inteligente cu senzori de presiune:} Instalează o barieră financiară și logistică prohibitivă, cerând modificări structurale ale locuinței și fiind susceptibile la zgomot acustic sau perturbări mecanice.
\end{enumerate}

În acest context, tehnologia emergentă de detecție ambientală pasivă bazată pe semnale de radiofrecvență (\textit{Device-Free Radio Frequency Sensing}), integrată în paradigma de comunicație și teledetecție sincronă (\textit{Integrated Sensing and Communication} -- ISAC) \cite{yang2022integrated}, oferă o alternativă complet non-intruzivă, privată și economic fezabilă. Semnalele Wi-Fi standard de 2.4 GHz pătrund prin pereți nemetalici, nu depind de iluminare, respectă în totalitate intimitatea utilizatorului (necaptând imagini vizuale) și exploatează infrastructura radio deja omniprezentă în spațiile de locuit.

\section{Obiectivele Tezei}
Obiectivul principal al acestei lucrări constă în proiectarea, implementarea fizică și validarea experimentală a unui sistem complet de monitorizare ambientală și detecție în timp real a căderilor umane, bazat pe extragerea și procesarea informației despre starea canalului radio (CSI -- \textit{Channel State Information}) pe o rețea triunghiulară de noduri cu cost redus bazate pe microcontrolere ESP32 și un analizor radio Ubertooth One.

Obiectivele specifice urmărite sunt:
\begin{enumerate}
    \item \textbf{Dezvoltarea unui firmware dedicat pentru extragerea CSI:} Configurarea modului promiscuu pe cipul ESP32 și a structurilor Wi-Fi CSI pe 52 de subpurtătoare active OFDM cu transmitere serială binară de mare viteză (921600 baud, 0xA55A magic header) la o cadență de eșantionare stabilă.
    \item \textbf{Proiectarea topologiei geometrice triunghiulare:} Amplasarea strategică a 3 noduri transmițător/receptor în încăpere pentru a forma o rețea confocală de elipsoide Fresnel (trei legături radio active: $L_{12}$, $L_{23}$, $L_{31}$), asigurând acoperire spațială 2D completă și redundanță la umbrire.
    \item \textbf{Dezvoltarea lanțului de procesare digitală a semnalului (DSP):} Implementarea unui pipeline analitic compus din filtrare Hampel (eliminare outlieri), filtrare trece-bandă Butterworth de ordinul 4 (0.2--12 Hz pentru izolarea dinamicii umane de derivă statică) și Analiza Componentelor Principale (PCA) pentru compresia dimensiunii spectrale pe prima componentă reprezentativă ($PC_1$).
    \item \textbf{Proiectarea unui automat cu stări finite analitic (FSM):} Dezvoltarea unui clasificator determinist care coroborează panta de colaps energetic simultan pe multiple legături radio, spectrul micro-Doppler ($>35\text{ Hz}$) și o fereastră obligatorie de imobilitate post-impact ($2.5\text{ s}$), garantând imunitate la alarme false cauzate de mersul rapid sau așezarea controlată.
    \item \textbf{Construirea unei aplicații desktop de monitorizare în timp real (PyQt6 + pyqtgraph):} Crearea unei interfețe grafice la 50--60 FPS ce randează poziția 2D a corpului în coordonate carteziene, deformarea dinamică a zonelor Fresnel, cascada spectrală waterfall și declanșează alerte de urgență vizuale și acustice.
    \item \textbf{Validarea riguroasă pe hardware fizic:} Colectarea și evaluarea datelor experimentale în trei scenarii controlate: Test A (Cameră goală), Test B (Mers normal) și Test C (Cădere reală simulată), determinând latența de decizie și acuratețea algoritmului.
\end{enumerate}

\section{Contribuții Originale}
Prezenta lucrare aduce următoarele contribuții originale:
\begin{itemize}
    \item \textbf{Arhitectură hardware replicabilă și accesibilă:} Spre deosebire de sistemele din literatura academică bazate pe plăci speciale Intel 5300 NIC \cite{halperin2011tool} sau stații SDR de laborator USRP \cite{adib2013see} care costă mii de euro și necesită cabluri coaxiale pentru sincronizarea ceasului, sistemul dezvoltat utilizează module accesibile ESP32 (sub 5 euro/bucată), demonstrând că detecția de cădere este posibilă pe platforme microcontroler embedded low-cost.
    \item \textbf{Criteriul de colaps simultan multi-link Fresnel:} Formularea unei condiții matematice robuste de impact bazate pe panta derivatei anvelopei $|dA/dt|$ pe cel puțin două legături simultan, combinată cu profilul Doppler STFT, eliminând necesitatea antrenării unor rețele neuronale de mari dimensiuni (deep learning) dependente de volume masive de date etichetate.
    \item \textbf{Algoritm hibrid de urmărire cartografică 2D:} Reconstrucția traiectoriei pasive prin optimizarea intersecției ecuațiilor eliptice Fresnel corespunzătoare celor 3 legături geometrice.
    \item \textbf{Suita software completă și deschisă:} Aplicație desktop integrată PyQt6 cu procesare asincronă pe fir de execuție separat (\texttt{QThread}), capabilă de streaming direct pe porturile seriale USB sau mod replay determinist.
\end{itemize}

\section{Structura Lucrării}
Lucrarea este organizată în șapte capitole, după cum urmează:
\begin{itemize}
    \item \textbf{Capitolul 1 (Introducere):} Prezintă contextul general, motivația socio-medicală, obiectivele tezei, contribuțiile originale și structura generală a documentului.
    \item \textbf{Capitolul 2 (Fundamente Teoretice și Stadiul Actual):} Expune teoria propagării radiofrecvenței, modelul zonelor Fresnel, diferențierea dintre RSSI și CSI, modulația OFDM conform standardului IEEE 802.11, precum și analiza comparativă a celor mai relevante lucrări din literatura de specialitate.
    \item \textbf{Capitolul 3 (Arhitectura Sistemului Hardware):} Detaliază configurația nodurilor ESP32, rolul analizorului Ubertooth One, protocolul de injectare a cadrelor și decodarea pachetelor de telemetrie binară pe interfața serială UART.
    \item \textbf{Capitolul 4 (Procesarea Semnalului și Algoritmi de Detecție):} Formalizează etapele matematice ale pipeline-ului DSP (filtrare outlieri, Butterworth, PCA, spectrogramă STFT Doppler) și logica de tranziție a automatului cu stări finite (FSM).
    \item \textbf{Capitolul 5 (Implementarea Software a Aplicației Desktop):} Descrie structura modulară a aplicației PyQt6, motorul de date \texttt{CSIDataEngine}, componentele grafice de randare (\texttt{RoomMapWidget}, \texttt{WaterfallWidget}, \texttt{DopplerWidget}, \texttt{TelemetryWidget}) și mecanismele de siguranță la erori.
    \item \textbf{Capitolul 6 (Rezultate Experimentale și Discuții):} Prezintă metodologia de testare pe bancul de lucru fizic, analizează comparativ seturile de date înregistrate (Test A, Test B, Test C), demonstrează comportamentul la impact și calculează latența de decizie.
    \item \textbf{Capitolul 7 (Concluzii și Direcții Viitoare):} Sintetizează realizările obținute, subliniază concluziile practice și schițează perspectivele de extindere a cercetării.
\end{itemize}
